\documentclass[10pt,a4paper]{amsart}
\usepackage[margin=19mm]{geometry}
\usepackage{amsmath,amssymb,mathtools}
\usepackage[scr=rsfs,cal=euler]{mathalfa}
\usepackage{xfrac}
\usepackage[unicode,hidelinks,bookmarks=false]{hyperref}
\newcommand{\ie}{i.e.\ }
\newcommand{\cf}{cf.\ }
\newcommand{\resp}{respectively\ }
\newcommand{\Subsection}{\subsection}
\providecommand{\nobreakdash}{\nobreak\hbox{-}}
\setlength{\emergencystretch}{2em}
\multlinegap0pt
\usepackage{amssymb}
\usepackage{mathtools}
\usepackage{xcolor}
\usepackage{bm}
\usepackage[shortlabels]{enumitem}
\usepackage{algorithm}\usepackage{algpseudocode}
\usepackage[normalem]{ulem}
\newtheorem{lemma}{Lemma}
\title{Spherical Harmonic Optimal Transport: Application to Climate Models Comparisons}
\author{Pierre Hou\'edry, Iskander Legheraba, L\'eo Buecher, Nicolas Courty}
\date{Source: arXiv:2605.18389v2 (CC BY 4.0)}
\begin{document}
\maketitle

\noindent\emph{Source and licence.} Spherical Harmonic Optimal Transport: Application to Climate Models Comparisons. Version arXiv:2605.18389v2 (CC BY 4.0), distributed by arXiv under the Creative Commons Attribution 4.0 International licence, http://creativecommons.org/licenses/by/4.0/, as recorded in the arXiv licence metadata for this exact version and shown on the abstract page https://arxiv.org/abs/2605.18389v2. Full licence notice: \texttt{licenses/arxiv-2605.18389-CC-BY-4.0.txt}.\par\smallskip

\noindent\textit{Editorial scope.} Counted unit: the article's lemma lem:block (source0.tex lines 673--694). The article's complete original proof of it is delivered with it (source0.tex lines 696--918, one \texttt{proof} environment of 223 lines). No mathematics is written, repaired or re-derived: the statement and the proof are the article's own lines, in the article's own notation, and no retained line is substituted or re-flowed.  No further source-local context is needed: every cross-reference of every retained line resolves inside the two delivered spans.  Nothing was omitted from any retained span, so the package carries no omission record.  No retained line contains a citation, so the wrapper carries no bibliography and no bibliography entry.  No retained line contains a figure environment, an image-inclusion command or drawing code, so no image is delivered and the package declares no asset dependency.  Editorial formatting only, in addition: an AMS \texttt{amsart} preamble that restates the article's own declarations and macros used by the retained lines, namely the environments \texttt{lemma} on separate counters, together with the standard packages those lines need.  The retained environments print in this document as Lemma 1 in order of appearance, whereas the article prints the same items with the numbering of its own full document.  The complete native source of the article as distributed in the arXiv e-print for this exact version is bundled unchanged as \texttt{sources/200853/source0.tex}.
\begin{lemma}[Block approximation]\label{lem:block}
Let \(X\) be a compact metric space, let \(P,Q \in \mathcal{R}_{+}(X)\), and let
$\pi \in \mathcal{R}_{+}(X^2)$
be such that $\mathrm{KL}(\pi^1\mid P)<+\infty$ and $\mathrm{KL}(\pi^2\mid Q)<+\infty.$
Let moreover $f:X^2 \rightarrow \mathbb{R}$ be a continuous function. Then for every \(\eta>0\), there exists $\pi_\eta \in \mathcal{R}_{+}(X^2)$
such that $\mathrm{KL}(\pi_\eta\mid P\otimes Q)<+\infty$, having the same mass as $\pi$ and satisfying
\begin{enumerate}
\item $\mathrm{KL}(\pi_\eta^1\mid P) \le \mathrm{KL}(\pi^1\mid P)$ and $\mathrm{KL}(\pi_\eta^2\mid Q)
\le
\mathrm{KL}(\pi^2\mid Q)$,
\item $
\left|
\int_{X^2} f\, d\pi
-
\int_{X^2} f\, \mathrm{d}\pi_\eta
\right|
\le \eta\, m(\pi).$
\end{enumerate}

Moreover, if \(m(P)=m(Q)\) and \(\pi\) is a coupling between \(P\) and \(Q\), then \(\pi_\eta\) is also a coupling between \(P\) and \(Q\).

\end{lemma}

\begin{proof}
Let \(\eta>0\). Since $f$ is continuous and \(X^2\) is compact, \(f\) is uniformly continuous. Hence there exists \(\delta>0\) such that
$d(x,x')\le \delta, ~ d(y,y') \le \delta \implies 
|f(x,y)-f(x',y')|
\le \eta.$

By compactness of \(X\), there exists a finite measurable partition
$\{A_1,\dots,A_N\}$
of \(X\) such that
$
\operatorname{diam}(A_i)\le \delta$ for all $i$.
Define
$I^+
:=
\{(i,j): P(A_i)>0,\ Q(A_j)>0\}$ 
and set
\[
\pi_\eta
:=
\sum_{(i,j)\in I^+}
\frac{\pi(A_i\times A_j)}{P(A_i)Q(A_j)}
(P\otimes Q)|_{A_i\times A_j}.
\]
By construction $\pi_\eta \ll P \otimes Q$ and $\mathrm{KL}(\pi_\eta\mid P\otimes Q)<+\infty$ and moreover
\[
\frac{d\pi_\eta}{d(P\otimes Q)}
=
\sum_{(i,j)\in I^+}
\frac{\pi(A_i\times A_j)}{P(A_i)Q(A_j)}
\mathbf 1_{A_i\times A_j}.
\]
We now study the marginals. Since
$\mathrm{KL}(\pi^1\mid P)<+\infty$
we have \(\pi^1\ll P\), thus the Radon-Nikodym derivative $\frac{d\pi^1}{dP}$
exists and for every measurable \(B\subset X\), one has
\[
\pi_\eta^1(B) =
\pi_\eta(B\times X) =
\sum_{(i,j)\in I^+}
\frac{\pi(A_i\times A_j)}{P(A_i)Q(A_j)}
\,P(B\cap A_i)\,Q(A_j)=
\sum_i
\frac{\pi^1(A_i)}{P(A_i)}
\,P(B\cap A_i),
\]
as \(\{A_j\}_j\) is a partition of \(X\). Therefore,
\[
\pi_\eta^1
=
\sum_i
\frac{\pi^1(A_i)}{P(A_i)}
\,P|_{A_i},
\]
and thus
\[
\frac{d\pi_\eta^1}{dP}
=
\sum_i
\left(
\frac1{P(A_i)}
\int_{A_i}\frac{d\pi^1}{dP}\, \mathrm{d}P
\right)\mathbf 1_{A_i}.
\]
Recalling that
\[
\mathrm{KL}(\mu\mid \nu)
=
\int
\varphi\!\left(\frac{d\mu}{d\nu}\right)\mathrm{d}\nu
\text{ with }
\varphi(t):= t\log t - t +1,
\]
the convexity of \(\varphi\) and Jensen's inequality in normalized form yield
\[
\varphi\!\left(
\frac1{P(A_i)}
\int_{A_i}\frac{d\pi^1}{dP}\, \mathrm{d}P
\right)
\le
\frac1{P(A_i)}
\int_{A_i}
\varphi \left(\frac{d\pi^1}{dP}\right) \mathrm{d}P.
\]
Multiplying by \(P(A_i)\) and summing over \(i\), we obtain
\[
\mathrm{KL}(\pi_\eta^1\mid P)
=
\sum_i
P(A_i)
\varphi\!\left(
\frac1{P(A_i)}
\int_{A_i}\frac{d\pi^1}{dP}\, \mathrm{d}P
\right) 
\le
\sum_i
\int_{A_i}
\varphi \left(\frac{d\pi^1}{dP}\right)\, \mathrm{d}P 
=
\mathrm{KL}(\pi^1\mid P).
\]
This proves the first part of (1). The proof of the second part is identical.

We now prove that $\pi$ and $\pi_\eta$ have the same mass. Observe that if \((i,j)\notin I^+\), then
\[
P(A_i)=0
\qquad
\text{or}
\qquad
Q(A_j)=0.
\]
If \(P(A_i)=0\), then since \(\pi^1\ll P\),
\[
\pi(A_i\times A_j)
\le
\pi(A_i\times X)
=
\pi^1(A_i)
=
0,
\]
as \(\pi^1\ll P\).
Similarly, if \(Q(A_j)=0\), then $\pi(A_i\times A_j)=0$.
Consequently $
\pi(A_i\times A_j)=0$
$(i,j)\notin I^+$.
Therefore,
\[
m(\pi_\eta)
=
\sum_{(i,j)\in I^+}
\frac{\pi(A_i\times A_j)}{P(A_i)Q(A_j)}
(P\otimes Q)(A_i\times A_j) 
=
\sum_{(i,j)\in I^+}
\pi(A_i\times A_j) 
=
\sum_{i,j}
\pi(A_i\times A_j) 
=
m(\pi).
\]

Finally, we prove (2). From what precedes, as $\pi(A_i\times A_j)=0$ for all $(i,j)\notin I^+$, one has
\[
\int_{X^2} f\, \mathrm{d}\pi = \sum_{(i,j)\in I^+}
\int_{A_i\times A_j} \int f \mathrm{d}\pi.
\]
Define
\[
\bar f_{ij}
:=
\frac1{P(A_i)Q(A_j)}
\int_{A_i\times A_j}
f\, \mathrm{d}(P\otimes Q).
\]
Then
\[
\int_{X^2} f\, \mathrm{d}\pi_\eta
=
\sum_{(i,j)\in I^+}
\pi(A_i\times A_j)\bar f_{ij},
\]
hence
\[
\int_{X^2} f\, \mathrm{d}\pi
-
\int_{X^2} f\, \mathrm{d}\pi_\eta
=
\sum_{(i,j)\in I^+}
\int_{A_i\times A_j}
(f-\bar f_{ij})\, \mathrm{d}\pi.
\]

Fix \((i,j)\in I^+\) and \((x,y)\in A_i\times A_j\). Then
\[
\begin{aligned}
|f(x,y)-\bar f_{ij}|
&=
\left|
\frac1{P(A_i)Q(A_j)}
\int_{A_i\times A_j}
(f(x,y)-f(x',y'))
\,\mathrm{d}(P\otimes Q)(x',y')
\right| \\
&\le
\frac1{P(A_i)Q(A_j)}
\int_{A_i\times A_j}
|f(x,y)-f(x',y')|
\,\mathrm{d}(P\otimes Q)(x',y').
\end{aligned}
\]
Since the $A_i$ have diameter less than $\delta$ and $(x,y) \in A_i \times A_j$, the uniform continuity property mentioned at the beginning gives, for all $(x,y) \in A_i \times A_j, ~
|f(x,y)-\bar f_{ij}|
\le \eta.
$
Consequently,
\[
\left|
\int_{X^2} f\, \mathrm{d}\pi
-
\int_{X^2} f\, \mathrm{d}\pi_\eta
\right|
\le
\sum_{(i,j)\in I^+}
\int_{A_i\times A_j}
|f-\bar f_{ij}|\, \mathrm{d}\pi 
\le
\eta
\sum_{(i,j)\in I^+}
\pi(A_i\times A_j) 
=
\eta\, m(\pi).
\]

Assume now that \(m(P)=m(Q)\) and that \(\pi\) is a coupling between \(P\) and \(Q\), that is $\pi^1 = P$ and $\pi^2 = Q$. Then $\frac{\pi^1(A_i)}{P(A_i)}=1$
thus $\pi_\eta^1
=
\sum_i P|_{A_i}
=
P$,
since \(\{A_i\}\) is a partition of \(X\). Similarly, $
\pi_\eta^2=Q$ and therefore \(\pi_\eta\) is also a coupling between \(P\) and \(Q\).
\end{proof}

\end{document}
